\chapter{فهرستِ کاملِ واژگانِ کلیدی}%
\label{app:keywords}
% longtable ships its pages from its own output routine, so the usual
% running header would keep the last mark from before the table.
\fancypagestyle{keywordsappendix}{%
  \fancyhf{}%
  \fancyhead[RO,LE]{\small\appendixname\ \thechapter.\ \ فهرستِ کاملِ واژگانِ کلیدی}%
  \fancyfoot[C]{\thepage}}
\pagestyle{keywordsappendix}

این پیوست مرجعِ کاملِ معادل‌های فارسی و انگلیسیِ زبانِ سلام است. هر واژه را
می‌توان به هر دو زبان نوشت و هر دو به یک برنامهٔ یکسان تبدیل می‌شوند؛ اما در یک
فایل باید یک زبان را نگه داشت. سلام زبانِ فایل را از واژه‌های کلیدی‌اش تشخیص
می‌دهد و پیام‌های خطا را هم به همان زبان می‌نویسد.

\section{واژگانِ کلیدی}

جدولِ زیر همهٔ واژه‌های کلیدیِ سلام را همراه با معادلِ انگلیسی‌شان نشان می‌دهد.

\begin{longtable}{l l p{0.5\linewidth}}
\toprule
\textbf{فارسی} & \textbf{انگلیسی} & \textbf{معنا} \\
\midrule
\endhead
\code{تابع} & \code{func} & تعریفِ تابع \\
\code{ریشه} & \code{main} & نامِ تابعِ ریشه (نقطهٔ آغازِ برنامه) \\
\code{برگشت} & \code{ret} & بازگرداندنِ مقدار \\
\code{اگر} & \code{if} & شرط \\
\code{وگرنه} & \code{else} & شاخهٔ جایگزین (با شرط بعد از آن، همان «وگرنه اگر») \\
\code{تا} & \code{until} & حلقه‌ای که \emph{تا وقتی} شرط برقرار است تکرار می‌شود \\
\code{تکرار} & \code{repeat} & حلقهٔ شمارشی \\
\code{تا} & \code{to} & کرانِ بالای بازه (درونِ \code{تکرار}) \\
\code{هر} & \code{each} & پیمایشِ عناصرِ یک مجموعه \\
\code{هر} & \code{by} & گامِ بازه (\code{تکرار ... تا ... هر ...}) \\
\code{در} & \code{in} & مجموعهٔ پیمایش‌شونده (با \code{هر}) و متغیرِ شمارنده (با \code{تکرار}) \\
\code{بشکن} & \code{break} & خروج از حلقه یا از یک حالتِ \code{ترابرد} \\
\code{گذر} & \code{continue} & پرش به تکرارِ بعد \\
\code{پایان} & \code{end} & پایانِ بلوک \\
\code{ناپایا} & \code{mut} & متغیرِ تغییرپذیر \\
\code{پایا} & \code{const} & ثابت \\
\code{گونه} & \code{type} & تعریفِ نامِ تازه برای یک گونه \\
\code{ساختار} & \code{struct} & ساختار \\
\code{جداشمار} & \code{enum} & شمارشی \\
\code{میانجی} & \code{interface} & میانجی \\
\code{کاربست} & \code{impl} & پیاده‌سازیِ یک میانجی برای یک گونه \\
\code{بر} & \code{on} & گونهٔ هدف (در \code{کاربست ... بر ...}) \\
\code{این} & \code{this} & نمونهٔ جاری درونِ متدها \\
\code{همخوان} & \code{match} & مقایسهٔ یک مقدار با چند الگو (بی‌ریزش) \\
\code{ترابرد} & \code{switch} & گزینش میانِ چند حالت (با ریزش به حالتِ بعد) \\
\code{دیرکن} & \code{defer} & اجرای به‌تعویق‌افتاده در پایانِ تابع \\
\code{واردسازی} & \code{import} & واردسازیِ بستهٔ کتابخانهٔ استاندارد \\
\code{فراخوانی} & \code{include} & آوردنِ فایلِ سلامِ خودتان \\
\code{بسته} & \code{package} & تعریفِ بسته \\
\code{همگانی} & \code{pub} & دسترسیِ همگانی \\
\code{درون‌داد} & \code{extern} & اعلانِ تابع‌های بیرونیِ \lr{C} \\
\code{برون‌داد} & \code{export} & بیرون‌دادنِ تابعِ سلام با نامِ ثابتِ \lr{C} \\
\code{برگردان} & \code{as} & تبدیلِ گونه \\
\code{درست} & \code{true} & مقدارِ درست \\
\code{نادرست} & \code{false} & مقدارِ نادرست \\
\code{پوچ} & \code{null} & اشاره‌گرِ تهی \\
\code{چاپ} & \code{print} & چاپ بدونِ خطِ تازه \\
\code{سرچاپ} & \code{println} & چاپ با خطِ تازه \\
\code{نادرست‌چاپ} & \code{printerr} & چاپ در خروجیِ خطا \\
\code{نادرست‌سرچاپ} & \code{printerrln} & چاپ در خروجیِ خطا با خطِ تازه \\
\code{ورودی} & \code{input} & خواندنِ یک خط از ورودی \\
\code{صفحه} & \code{page} & صفحهٔ وب \\
\code{بخش} & \code{component} & بخشِ بازکاربردیِ صفحه \\
\code{و} & \code{and} & «و»ی منطقی \\
\code{یا} & \code{or} & «یا»ی منطقی \\
\code{وارونه} & \code{not} & نقیضِ منطقی \\
\code{برابر} & \code{eq} & برابری (همان \code{==}) \\
\code{نابرابر} & \code{neq} & نابرابری (همان \code{!=}) \\
\code{کوچکتر} & \code{lt} & کوچک‌تر (همان \code{<}) \\
\code{بزرگتر} & \code{gt} & بزرگ‌تر (همان \code{>}) \\
\code{کوچکتربرابر} & \code{lte} & کوچک‌تر یا برابر (همان \code{<=}) \\
\code{بزرگتربرابر} & \code{gte} & بزرگ‌تر یا برابر (همان \code{>=}) \\
\code{درخط} & \code{inline} & پیشنهادِ درج‌شدنِ تابع در محلِ فراخوانی \\
\code{نادرخط} & \code{noinline} & منعِ درج‌شدنِ تابع \\
\code{ناب} & \code{pure} & تابعِ بدونِ اثرِ جانبی \\
\code{نابرگشت} & \code{noret} & تابعی که هرگز برنمی‌گردد \\
\code{بی‌کاره} & \code{deprecated} & تابعِ منسوخ \\
\bottomrule
\end{longtable}

\Needspace{14\baselineskip}
\section{واژه‌های وابسته به جایگاه}

این واژه‌ها فقط در جایگاهِ ویژهٔ خود کلیدی‌اند و در جاهای دیگر می‌توانند نام
باشند:

\begin{longtable}{l l p{0.5\linewidth}}
\toprule
\textbf{فارسی} & \textbf{انگلیسی} & \textbf{کاربرد} \\
\midrule
\endhead
\code{پیوند} & \code{link} & پیوند به کتابخانهٔ بیرونی (\code{پیوند پویا "..."}) \\
\code{پویا} & \code{dynamic} / \code{dyn} & پیوندِ پویا؛ و پیش از نامِ میانجی، گونهٔ پویا \\
\code{ایستا} & \code{static} & پیوندِ ایستا؛ و متدِ ایستای ساختار \\
\code{چارچوب} & \code{framework} & پیوند به چارچوبِ مک \\
\code{شامل} & \code{use} & جاسازیِ یک ساختار در ساختارِ دیگر \\
\code{تعریف‌شده} & \code{defined} & بررسیِ وجودِ یک پرچمِ زمانِ کامپایل \\
\bottomrule
\end{longtable}

\section{تابع‌های درونی}

این تابع‌ها بخشی از خودِ زبان‌اند و به واردسازی نیاز ندارند. در فایلِ فارسی
باید نامِ فارسی‌شان را نوشت؛ نامِ انگلیسی در آن‌جا خطاست.

\begin{longtable}{l l p{0.5\linewidth}}
\toprule
\textbf{فارسی} & \textbf{انگلیسی} & \textbf{کار} \\
\midrule
\endhead
\code{طول} & \code{len} & طولِ آرایه، رشته یا بردار \\
\code{اندازه‌گونه} & \code{sizeof} & اندازهٔ یک گونه به بایت \\
\code{آرگومان‌ها} & \code{args} & آرگومان‌های خطِ فرمان \\
\code{متغیر‌محیطی} & \code{env} & خواندنِ متغیرِ محیطی \\
\code{زبان} & \code{lang} & زبانِ فایلِ برنامه \\
\code{بازکردن} & \code{open} & بازکردنِ پرونده \\
\code{فهرست‌پوشه} & \code{listdir} & فهرستِ پرونده‌های یک پوشه \\
\code{درهم} & \code{hash} & درهم‌سازیِ یک مقدار \\
\code{کدنویسه} & \code{char\_code} & کدِ عددیِ یک نویسه \\
\code{نویسه‌ازکد} & \code{char\_from\_code} & نویسه از روی کدِ عددی \\
\code{مقایسه‌رشته} & \code{strcmp} & مقایسهٔ دو رشته \\
\code{نخ‌ساز} & \code{spawn} & ساختنِ نخِ تازه \\
\code{نخ‌پیوند} & \code{join} & انتظار برای پایانِ نخ \\
\code{فراخوان‌دستگیره} & \code{callhandler} & صدا زدنِ یک دستگیرهٔ تابع \\
\code{اتمی‌بخوان} & \code{atomic\_load} & خواندنِ اتمی \\
\code{اتمی‌بنویس} & \code{atomic\_store} & نوشتنِ اتمی \\
\code{اتمی‌بیفزا} & \code{atomic\_add} & افزایشِ اتمی \\
\code{اتمی‌جابجا} & \code{atomic\_swap} & جابه‌جاییِ اتمی \\
\code{اتمی‌مقایسه} & \code{atomic\_cas} & مقایسه و جابه‌جاییِ اتمی \\
\bottomrule
\end{longtable}

\section{گونه‌های داده}

جدولِ زیر نامِ فارسی و انگلیسیِ گونه‌های داخلیِ زبان را نشان می‌دهد. رقم‌های
نامِ گونه را می‌توان فارسی یا لاتین نوشت (\code{صحیح۶۴} همان \code{صحیح64} است).

\begin{longtable}{l l p{0.36\linewidth}}
\toprule
\textbf{فارسی} & \textbf{انگلیسی} & \textbf{توضیح} \\
\midrule
\endhead
\code{صحیح} / \code{صحیح۳۲} & \code{int} / \code{i32} & عددِ صحیحِ ۳۲ بیتی \\
\code{صحیح۸} & \code{i8} & عددِ صحیحِ ۸ بیتی \\
\code{صحیح۱۶} & \code{i16} & عددِ صحیحِ ۱۶ بیتی \\
\code{صحیح۶۴} & \code{i64} & عددِ صحیحِ ۶۴ بیتی \\
\code{طبیعی} / \code{طبیعی۳۲} & \code{uint} / \code{u32} & عددِ بی‌علامتِ ۳۲ بیتی \\
\code{طبیعی۸} & \code{u8} & عددِ بی‌علامتِ ۸ بیتی \\
\code{طبیعی۱۶} & \code{u16} & عددِ بی‌علامتِ ۱۶ بیتی \\
\code{طبیعی۶۴} & \code{u64} & عددِ بی‌علامتِ ۶۴ بیتی \\
\code{اندازه} & \code{size} & اندازه (هم‌اندازهٔ اشاره‌گر) \\
\code{اندازه مثبت} & \code{usize} & اندازهٔ بی‌علامت \\
\code{اعشار} / \code{اعشار۳۲} & \code{float} / \code{f32} & عددِ اعشاریِ ۳۲ بیتی \\
\code{اعشار۶۴} & \code{f64} & عددِ اعشاریِ ۶۴ بیتی (گونهٔ هر عددِ اعشاریِ نوشته‌شده) \\
\code{منطقی} & \code{bool} & درست/نادرست \\
\code{نویسه} & \code{char} & یک بایت (نویسهٔ \lr{ASCII}) \\
\code{یونیکد} & \code{uchar} & یک نویسهٔ یونیکد \\
\code{رشته} & \code{str} & متن \\
\code{تهی} & \code{void} & بی‌مقدار \\
\code{وکتور<T>} & \code{Vector<T>} & بردارِ رشدپذیر \\
\code{نگاشت<K، V>} & \code{HashMap<K, V>} & نگاشتِ کلید به مقدار \\
\code{پیمایشگرنگاشت} & \code{MapIter} & پیمایشگرِ نگاشت \\
\code{پرونده} & \code{File} & پروندهٔ باز \\
\code{گوناگون<...>} & \code{Variant<...>} & مقداری از یکی از چند گونه \\
\bottomrule
\end{longtable}

\section{بسته‌های کتابخانهٔ استاندارد}

پرکاربردترین بسته‌های کتابخانهٔ استاندارد در جدولِ زیر آمده‌اند. در فایلِ
انگلیسی با نامِ انگلیسی (\code{import math}) و در فایلِ فارسی با نامِ فارسی
(\code{واردسازی ریاضی}) واردشان کنید؛ بسته‌های تودرتو با نقطه نوشته می‌شوند.

\begin{longtable}{l l l}
\toprule
\textbf{واردسازی} & \textbf{نامِ فارسی} & \textbf{کاربرد} \\
\midrule
\endhead
\code{math} & \code{ریاضی} & تابع‌های ریاضی \\
\code{str} & \code{رشته} & پردازشِ متن \\
\code{io} & \code{ورودی خروجی} & ورودی/خروجی و پرونده \\
\code{os} & \code{سیستم عامل} & سیستم‌عامل، آرگومان‌ها، محیط \\
\code{time} & \code{زمان} & زمان و تاریخ \\
\code{rand} & \code{تصادفی} & عددهای تصادفی \\
\code{fmt} & \code{قالب بندی} & قالب‌بندیِ متن \\
\code{conv} & \code{تبدیل} & تبدیلِ عدد و متن \\
\code{sort} & \code{مرتب سازی} & مرتب‌سازی \\
\code{collections} & \code{مجموعه ها} & ساختارهای داده \\
\code{encoding.json} & \code{کدگذاری.جیسون} & دادهٔ \lr{JSON} \\
\code{sync} & \code{همگام سازی} & قفل‌ها و گروهِ انتظار \\
\code{net.http} & \code{شبکه.اچ تی تی پی} & درخواست‌ها و کارسازهای وب \\
\code{testing} & \code{آزمایش} & آزمون‌نویسی \\
\code{mem} & \code{حافظه} & مدیریتِ حافظه \\
\code{option} & \code{اختیاری} & مقدارهای اختیاری \\
\code{regex} & \code{عبارت باقاعده} & عبارت‌های باقاعده \\
\code{db.sqlite} & \code{دیتابیس.اس کیو لایت} & پایگاه‌دادهٔ \lr{SQLite} \\
\code{db.mysql} & \code{دیتابیس.مای اس کیو ال} & پایگاه‌دادهٔ \lr{MySQL} \\
\bottomrule
\end{longtable}

\section{فرمان‌های کامپایلر}

فرمان‌های پرکاربردِ ابزارِ \code{salam} در جدولِ زیر آمده‌اند؛ فهرستِ کامل را
\code{salam help} نشان می‌دهد.

\begin{longtable}{p{0.4\linewidth} p{0.5\linewidth}}
\toprule
\textbf{فرمان} & \textbf{کار} \\
\midrule
\endhead
\code{salam run app.salam} & ساخت و اجرا بدونِ نگه‌داشتنِ فایلِ اجرایی \\
\code{salam exec app.salam} & اجرای بی‌درنگ با مفسر \\
\code{salam check app.salam} & وارسیِ خطاها بدونِ ساخت \\
\code{salam build app.salam --output=app} & ساختِ فایلِ اجرایی \\
\code{salam build app.salam --release} & ساختِ بهینه بدونِ وارسیِ مرزها \\
\code{salam build app.salam -DFLAG} & تعریفِ پرچمِ زمانِ کامپایل \\
\code{salam build app.salam --emit-c} & نوشتنِ کدِ \lr{C} تولیدشده \\
\code{salam obj app.salam} & ساختِ فایلِ شیء (\lr{.o}) \\
\code{salam page build page.salam} & ساختِ صفحهٔ وب (\lr{HTML/CSS/JS}) \\
\code{salam page build page.salam --inline} & صفحهٔ یکپارچه \\
\code{salam format app.salam} & قالب‌بندیِ کد \\
\code{salam format --check} & گزارشِ فایل‌های نامرتب \\
\code{salam translate fa app.salam} & برگرداندنِ کدِ انگلیسی به فارسی \\
\code{salam new myapp --lang=fa} & ساختنِ پروژهٔ تازه \\
\code{salam debug app.salam} & ساخت و اجرای اشکال‌زدا \\
\code{salam memcheck app.salam} & وارسیِ نشتِ حافظه \\
\code{salam inspect app.salam --emit-tokens} & نمایشِ نشانه‌ها \\
\code{salam inspect app.salam --emit-ast} & نمایشِ درختِ نحوی \\
\bottomrule
\end{longtable}

\section{راهنماهای فایل}

این راهنماها درونِ خودِ کد نوشته می‌شوند و رفتارِ کامپایلر را تعیین می‌کنند.

\begin{longtable}{l p{0.6\linewidth}}
\toprule
\textbf{راهنما} & \textbf{معنا} \\
\midrule
\endhead
\code{@en "..."} & نامِ انگلیسیِ یک تعریف \\
\code{@fa "..."} & نامِ فارسیِ یک تعریف \\
\code{اگر / وگرنه / پایان} در سطحِ بالا & شرطِ زمانِ کامپایل (همان کلیدواژه‌های شرطِ معمولی) \\
\code{سلام سیستم ویندوز} و مانندِ آن & پرچم‌های سیستم‌عامل (\code{SALAM\_OS\_WINDOWS}، \ldots) \\
\code{تعریف‌شده(FLAG)} & بررسیِ وجودِ یک پرچمِ زمانِ کامپایل (\code{defined}) \\
\bottomrule
\end{longtable}

\clearpage
\pagestyle{bookpages}
